\chapter{Reaction Mechanisms of Complexes}\label{ch:b3:complex-mechanisms}

Dissolve a chromium(III) salt in water labelled with oxygen-18, and days pass
before the labelled water has replaced the water molecules bound to the
metal; with a copper(II) salt the exchange is complete within the time of
mixing. The reaction is the same, one water molecule leaving a metal ion and
another taking its place, yet its rate varies enormously from one ion to the
next, and the number of $d$ electrons predicts much of the variation. The way
ligands are replaced decides how platinum anticancer drugs are made and how
they act; the way electrons move between metal ions decides the speed of
respiration and of every battery. This chapter treats the two families of
reactions of complexes: substitution and electron transfer, the latter with
the theory of Rudolph Marcus.

\begin{recall}
The Year 2 volume described complexes, the denticity of ligands, bridging
ligands, $fac$ and $mer$ isomers, ligand exchange as an elementary step, and
crystal-field stabilisation energies with high- and low-spin configurations;
the Year 1 volume, mechanisms with their steady-state and rate-determining-step
approximations. \Cref{ch:b3:rate-theories} gave the \omterm{thm:b3:rate-theories:eyring}{Eyring equation} and the
meaning of $\Delta^{\ddagger}S^\circ$, \cref{ch:b3:electrode-kinetics} electron transfer at an
electrode, and \cref{ch:b3:complex-spectra-magnetism} the \omterm{def:b3:complex-spectra-magnetism:ligand-field-term}{ligand-field terms} and the
Jahn--Teller effect.
\end{recall}

\section{Lability and inertness}

\begin{definition}[Labile and inert complexes]\label{def:b3:complex-mechanisms:labile}
A complex is \emph{labile}\index{labile complex} if its ligands are replaced
quickly, and a \emph{kinetically inert complex}\index{kinetically inert complex}
if they are replaced slowly; following Taube, a complex is called labile when
its reactions with a ligand at \qty{0.1}{mol/L} and \qty{25}{\celsius} are over within about a
minute. Lability is a kinetic property, unrelated to the thermodynamic stability
of the complex.
\end{definition}

Water exchange, $\mathrm{[M(H_2O)_6]^{n+} + H_2O^{\ast} \to [M(H_2O)_5(H_2O^{\ast})]^{n+} + H_2O}$, is the simplest
substitution and sets a scale. Ions of the main groups exchange faster when they
are larger and less charged. Among transition-metal ions, those with $d^3$
(\ce{Cr^{3+}}) or low-spin $d^6$ (\ce{Co^{3+}}, \ce{Rh^{3+}}, \ce{Ir^{3+}}) configurations are inert, and so are
most second- and third-row metals, while $d^9$ (\ce{Cu^{2+}}) and high-spin $d^4$ (\ce{Cr^{2+}}),
distorted by the Jahn--Teller effect, are extremely \omterm{def:b3:complex-mechanisms:labile}{labile}.

\begin{proposition}[Ligand-field activation energy]\label{prop:b3:complex-mechanisms:lfae}
If a substitution goes through a five-coordinate square-pyramidal transition
state, the loss of ligand-field stabilisation on going there (the ligand-field
activation energy) is, in the simplest $\sigma$-only model, largest for low-spin $d^6$
($4\,Dq$) and next for $d^3$ and $d^8$ ($2\,Dq$), zero for $d^0$, high-spin $d^5$ and $d^{10}$, and negative
for high-spin $d^4$ and $d^9$.
\end{proposition}

\begin{proof}
Model each ligand as raising a $d$ orbital by $e_\sigma$ times the square of the orbital's
angular amplitude in its direction, normalised to 1 on the lobe axis: an axial
ligand gives $d_{z^2}$ 1, an equatorial one gives $d_{z^2}$ $\frac14$ and $d_{x^2-y^2}$ $\frac34$; the $t_{2g}$ orbitals point
between the ligands and get nothing. The octahedron gives $d_{z^2}$ and $d_{x^2-y^2}$ $3e_\sigma$ each
($\Delta_{\mathrm o} = 3e_\sigma = 10Dq$); the square pyramid, without one axial ligand, gives $d_{z^2}$ $2e_\sigma$ and
$d_{x^2-y^2}$ $3e_\sigma$. The stabilisation of $n$ electrons is their energy minus $n/5$ times the sum of
the five levels (the barycentre). For low-spin $d^6$ ($t_{2g}^6$): $0 - \frac65 \times 6e_\sigma$ in the
octahedron, $0 - \frac65 \times 5e_\sigma$ in the pyramid: a loss of $1.2e_\sigma = 4Dq$. For $d^3$: $\frac35e_\sigma =
2Dq$; for $d^8$ ($t_{2g}^6e_g^2$): $(5 - \frac85 \times 5) - (6 - \frac85 \times 6) = 0.6e_\sigma$. For $d^9$: $(7 - 9) - (9 -
10.8) = -0.2e_\sigma$.
\end{proof}

\begin{omfigure}
\begin{tikzpicture}
\begin{axis}[width=0.72\linewidth, height=5.4cm, xmin=-0.6, xmax=10.6, ymin=-1.2, ymax=4.6,
  axis x line=bottom, axis y line=left, ycomb, xtick={0,1,2,3,4,5,6,7,8,9,10},
  xlabel={$n$ ($d^n$)}, ylabel={activation energy / $Dq$}, tick label style={font=\small},
  label style={font=\small}]
\draw[black!40] (axis cs:-0.6,0) -- (axis cs:10.6,0);
\addplot[omDef, line width=4pt, xshift=-2.5pt] table[x=n, y=hs] {figdata/out/bachelor-3-ligand-field-activation.dat};
\addplot[omThm, line width=4pt, xshift=2.5pt, unbounded coords=discard] table[x=n, y=ls] {figdata/out/bachelor-3-ligand-field-activation.dat};
\end{axis}
\end{tikzpicture}

{\small Ligand-field activation energies for a dissociative path through a
square pyramid, in the $\sigma$-only model of the proposition (blue: high spin;
red: low spin). The largest values,
for low-spin $d^6$, $d^3$ and $d^8$, match the inert or slow ions (\ce{Co^{3+}}, \ce{Cr^{3+}}, \ce{Ni^{2+}});
the negative ones, for $d^4$ and $d^9$, the very \omterm{def:b3:complex-mechanisms:labile}{labile} Jahn--Teller ions.}
\end{omfigure}

\begin{definition}[Volume of activation]\label{def:b3:complex-mechanisms:activation-volume}
The \emph{volume of activation}\index{volume of activation} $\Delta V^{\ddagger}$ of an elementary
step is the difference between the partial molar volume of the \omterm{def:b3:rate-theories:activated-complex}{activated
complex} and those of the reactants.
\end{definition}

\begin{proposition}[Pressure dependence of a rate constant]\label{prop:b3:complex-mechanisms:pressure}
At constant temperature, $\displaystyle\Bigl(\frac{\partial\ln k}{\partial p}\Bigr)_T = -\frac{\Delta V^{\ddagger}}{RT}$.
\end{proposition}

\begin{proof}
By the \omterm{thm:b3:rate-theories:eyring}{Eyring equation} $\ln k = \text{const} - \Delta^{\ddagger}G/RT$, and $(\partial\Delta^{\ddagger}G/\partial p)_T = \Delta V^{\ddagger}$, as for
any Gibbs energy.
\end{proof}

A step in which a ligand leaves has $\Delta V^{\ddagger} > 0$; one in which a ligand enters has $\Delta V^{\ddagger}
< 0$. Measured at a few hundred megapascals, the activation volume is the most
direct diagnosis of a substitution mechanism, less blurred by solvation than
the \omterm{def:b3:rate-theories:activation}{entropy of activation}.

\section{Substitution mechanisms}

\begin{definition}[Dissociative, associative and interchange mechanisms]\label{def:b3:complex-mechanisms:mechanisms}
In a \emph{dissociative mechanism}\index{dissociative mechanism} (D) the leaving
ligand departs first, giving an intermediate of lower coordination number; in an
\emph{associative mechanism}\index{associative mechanism} (A) the entering ligand
binds first, giving an intermediate of higher coordination number; in an
\emph{interchange mechanism}\index{interchange mechanism} (I) the two ligands
exchange in a single step, with a bond-breaking ($I_d$) or bond-making ($I_a$)
character.
\end{definition}

\begin{proposition}[Rate laws of D and A mechanisms]\label{prop:b3:complex-mechanisms:d-a-laws}
For $\mathrm{ML_5X + Y \to ML_5Y + X}$: a D mechanism ($\mathrm{ML_5X} \rightleftharpoons \mathrm{ML_5} + \mathrm X$, $k_1$, $k_{-1}$; $\mathrm{ML_5} + \mathrm Y
\to \mathrm{ML_5Y}$, $k_2$) gives
\[
  v = \frac{k_1k_2[\mathrm{ML_5X}][\mathrm Y]}{k_{-1}[\mathrm X] + k_2[\mathrm Y]},
\]
first order in the complex and independent of $[\mathrm Y]$ at high $[\mathrm Y]$; an A mechanism
through a seven-coordinate intermediate in a steady state gives $v =
k[\mathrm{ML_5X}][\mathrm Y]$, first order in each.
\end{proposition}

\begin{proof}
D: steady state for $\mathrm{ML_5}$, $k_1[\mathrm{ML_5X}] = (k_{-1}[\mathrm X] + k_2[\mathrm Y])[\mathrm{ML_5}]$, and $v = k_2[\mathrm{ML_5}][\mathrm Y]$.
At high $[\mathrm Y]$, $v \to k_1[\mathrm{ML_5X}]$. A: $\mathrm{ML_5X} + \mathrm Y \rightleftharpoons \mathrm{ML_5XY}$ ($k_a$, $k_{-a}$), $\mathrm{ML_5XY} \to
\mathrm{ML_5Y} + \mathrm X$ ($k_b$); the steady state gives $v = \frac{k_ak_b}{k_{-a} + k_b}[\mathrm{ML_5X}][\mathrm Y]$.
\end{proof}

In water, the entering ligand is usually present at much lower concentration
than the solvent, and the true first step is the meeting of the complex and the
ligand.

\begin{definition}[Eigen--Wilkins mechanism]\label{def:b3:complex-mechanisms:eigen-wilkins}
The \emph{Eigen--Wilkins mechanism}\index{Eigen--Wilkins mechanism} of
substitution of an aqua complex is a fast pre-equilibrium forming an
\emph{outer-sphere complex}\index{outer-sphere complex}, in which the entering
ligand sits next to the complex, outside its coordination sphere, followed by
the rate-determining interchange of a water molecule for that ligand.
\end{definition}

\begin{theorem}[Eigen--Wilkins rate law]\label{thm:b3:complex-mechanisms:eigen-wilkins}
With the outer-sphere equilibrium constant $K_{\mathrm{os}}$ and the interchange rate constant
$k_i$, the observed pseudo-first-order rate constant at excess $[\mathrm Y]$ is
\[
  k_{\mathrm{obs}} = \frac{K_{\mathrm{os}}k_i[\mathrm Y]}{1 + K_{\mathrm{os}}[\mathrm Y]} .
\]
\end{theorem}

\begin{proof}
The complex is distributed between free and outer-sphere forms, $[\mathrm{M\cdot Y}] = K_{\mathrm{os}}[\mathrm M][\mathrm Y]$,
with total $[\mathrm M]_{\mathrm{tot}} = [\mathrm M](1 + K_{\mathrm{os}}[\mathrm Y])$; the rate is $k_i[\mathrm{M\cdot Y}] = k_iK_{\mathrm{os}}[\mathrm Y][\mathrm M]_{\mathrm{tot}}/(1 + K_{\mathrm{os}}[\mathrm Y])$.
\end{proof}

At low $[\mathrm Y]$ the reaction is second order with $k = K_{\mathrm{os}}k_i$; $K_{\mathrm{os}}$, estimated from the
charges and the distance of closest approach, is nearly the same for all
ligands of the same charge, so that $k_i$, the water-exchange step, controls the
rate: substitution on an aqua ion is nearly independent of the entering
ligand, the signature of a dissociative interchange.

\begin{definition}[Conjugate-base mechanism]\label{def:b3:complex-mechanisms:conjugate-base}
The \emph{conjugate-base mechanism}\index{conjugate-base mechanism} of base
hydrolysis of an ammine complex, $\mathrm{[Co(NH_3)_5X]^{2+} + OH^- \to [Co(NH_3)_5(OH)]^{2+} + X^-}$, is
the deprotonation of an ammine ligand in a fast pre-equilibrium, followed by the
dissociative loss of $\mathrm X^-$ from the amido conjugate base, whose $\mathrm{NH_2^-}$ ligand
labilises the complex.
\end{definition}

\begin{proposition}[Rate law of base hydrolysis]\label{prop:b3:complex-mechanisms:conjugate-base}
With the deprotonation constant $K$ and the dissociation rate constant $k$ of the
conjugate base, $v = \dfrac{kK[\text{complex}]_{\mathrm{tot}}[\ce{OH-}]}{1 + K[\ce{OH-}]}$, second order at low $[\ce{OH-}]$.
\end{proposition}

\begin{proof}
As for the Eigen--Wilkins law: the complex is distributed between its acid and
conjugate-base forms in the ratio $1 : K[\ce{OH-}]$, and only the latter reacts.
\end{proof}

The rate law alone does not distinguish this mechanism from a direct attack of
\ce{OH-}; the evidence is that base hydrolysis needs an N--H proton (complexes without
one do not show it) and that the ammine protons exchange with the solvent much
faster than the hydrolysis.

\begin{method}[Diagnosing a substitution mechanism]\label{met:b3:complex-mechanisms:diagnose}
\begin{enumerate}
  \item Entering-group dependence: rate independent of Y (or saturating),
        dissociative; rate proportional to $[\mathrm Y]$ and sensitive to the nature of Y,
        associative.
  \item Activation volume: positive, dissociative; negative, associative.
  \item \omterm{def:b3:rate-theories:activation}{Entropy of activation}: positive for D, negative for A (with caution, since
        solvation contributes).
  \item Leaving-group dependence: a rate that follows the strength of the M--X bond
        points to bond breaking in the transition state.
\end{enumerate}
\end{method}

\begin{omfigure}
\begin{tikzpicture}
\begin{axis}[name=pd, omprofile, width=0.46\linewidth, height=4.6cm, xmin=0, xmax=10, ymin=0, ymax=10,
  xlabel={reaction coordinate}, ylabel={energy}, title={\small D: $\mathrm{ML_5X} \to \mathrm{ML_5} + \mathrm X \to \mathrm{ML_5Y}$}]
\addplot[omDef, very thick, smooth] coordinates {(0,2) (1.5,2.1) (3,7.5) (4.2,5.6) (5.0,5.4) (5.8,5.6) (7,7.0) (8.5,1.6) (10,1.5)};
\node[font=\scriptsize, anchor=north, align=center] at (axis cs:5,5.2) {ML$_5$ (CN 5)};
\end{axis}
\begin{axis}[at={(pd.east)}, anchor=west, xshift=1.0cm, omprofile, width=0.46\linewidth, height=4.6cm,
  xmin=0, xmax=10, ymin=0, ymax=10, xlabel={reaction coordinate}, ylabel={energy},
  title={\small A: $\mathrm{ML_5X} + \mathrm Y \to \mathrm{ML_5XY} \to \mathrm{ML_5Y}$}]
\addplot[omThm, very thick, smooth] coordinates {(0,2) (1.5,2.1) (3,6.5) (4.2,4.4) (5.0,4.2) (5.8,4.4) (7,7.2) (8.5,1.6) (10,1.5)};
\node[font=\scriptsize, anchor=north, align=center] at (axis cs:5,4.0) {ML$_5$XY (CN 7)};
\end{axis}
\end{tikzpicture}

{\small Energy profiles of a dissociative and an associative substitution
(schematic): both pass through an intermediate, of lower coordination number
(CN) in D and of higher coordination number in A. In the \omterm{def:b3:complex-mechanisms:mechanisms}{interchange mechanisms}
the well disappears and a single transition state remains.}
\end{omfigure}

An octahedral complex that substitutes through a square pyramid generally
keeps its configuration ($cis$ stays $cis$); one that passes through a trigonal
bipyramid, as many cobalt(III) complexes do in base hydrolysis, can give a
mixture of $cis$ and $trans$ products.

\section{Square-planar substitution and the trans effect}

Square-planar complexes of $d^8$ ions (\ce{Pt^{2+}}, \ce{Pd^{2+}}, \ce{Au^{3+}}, \ce{Ni^{2+}} with strong ligands)
are coordinatively unsaturated: an entering ligand can approach along the
empty axial direction. Their substitutions are associative, through a
five-coordinate trigonal-bipyramidal transition state.

\begin{proposition}[Two-term rate law]\label{prop:b3:complex-mechanisms:square-planar}
The substitution $\mathrm{[PtL_3X] + Y \to [PtL_3Y] + X}$ in a coordinating solvent S follows
$k_{\mathrm{obs}} = k_1 + k_2[\mathrm Y]$ at excess Y.
\end{proposition}

\begin{proof}
Two parallel associative paths: direct attack of Y ($k_2[\mathrm Y]$), and attack of the
solvent, present in large excess (pseudo-first-order $k_1$), giving $[\mathrm{PtL_3S}]$, which
is then rapidly converted by Y. Parallel first-order paths add.
\end{proof}

\begin{definition}[Trans effect, trans influence]\label{def:b3:complex-mechanisms:trans-effect}
The \emph{trans effect}\index{trans effect} is the acceleration of the substitution
of a ligand by the ligand trans to it: a kinetic effect, on the transition state.
The \emph{trans influence}\index{trans influence} is the weakening, in the ground
state, of the bond trans to a ligand, seen in bond lengths and vibrational
frequencies: a thermodynamic effect.
\end{definition}

The \omterm{def:b3:complex-mechanisms:trans-effect}{trans effect} follows the order $\ce{CO}$, $\ce{CN-}$, $\ce{C2H4} > \ce{PR3}$, $\ce{H-} > \ce{CH3-} > \ce{I-}$, $\ce{SCN-} >
\ce{Br-} > \ce{Cl-} > \ce{NH3}$, $\ce{H2O}$. Strong $\sigma$ donors weaken the trans bond (\omterm{def:b3:complex-mechanisms:trans-effect}{trans influence});
strong $\pi$ acceptors stabilise the five-coordinate transition state by taking
\omterm{def:b3:computational-chemistry:dft}{electron density} from the metal. Both raise the rate.

\begin{method}[Planning the synthesis of platinum(II) isomers]\label{met:b3:complex-mechanisms:pt-synthesis}
\begin{enumerate}
  \item At each step, the ligand replaced is the one trans to the ligand of
        highest \omterm{def:b3:complex-mechanisms:trans-effect}{trans effect} present.
  \item Order the additions so that this rule leads to the wanted isomer.
\end{enumerate}
\end{method}

\begin{omfigure}
\setchemfig{atom sep=2.2em}
\schemestart
\chemfig{Cl-Pt(-[2]Cl)(-[6]Cl)-Cl}
\arrow{->[\ce{NH3}]}
\chemfig{Cl-Pt(-[2]Cl)(-[6]Cl)-NH_3}
\arrow{->[\ce{NH3}]}
\chemfig{Cl-Pt(-[2]NH_3)(-[6]Cl)-NH_3}
\schemestop

\bigskip
\schemestart
\chemfig{H_3N-Pt(-[2]NH_3)(-[6]NH_3)-NH_3}
\arrow{->[\ce{Cl-}]}
\chemfig{H_3N-Pt(-[2]NH_3)(-[6]NH_3)-Cl}
\arrow{->[\ce{Cl-}]}
\chemfig{Cl-Pt(-[2]NH_3)(-[6]NH_3)-Cl}
\schemestop

{\small The \omterm{def:b3:complex-mechanisms:trans-effect}{trans effect} at work (charges omitted). Top: from $\ce{[PtCl4]^{2-}}$, the second
ammonia replaces a chloride trans to chloride (Cl$^-$ has the larger \omterm{def:b3:complex-mechanisms:trans-effect}{trans effect}),
giving the $cis$ isomer, cisplatin. Bottom: from $\ce{[Pt(NH3)4]^{2+}}$, the second chloride
replaces the ammonia trans to the first chloride, giving the $trans$ isomer.}
\end{omfigure}

\section{Electron transfer}

\begin{definition}[Outer- and inner-sphere electron transfer]\label{def:b3:complex-mechanisms:electron-transfer}
In \emph{outer-sphere electron transfer}\index{outer-sphere electron transfer} the
electron passes between two complexes whose coordination spheres stay intact;
in \emph{inner-sphere electron transfer}\index{inner-sphere electron transfer} it
passes through a ligand bridging the two metals in a transient binuclear
complex. A \emph{self-exchange reaction}\index{self-exchange reaction} is an
electron transfer between the two oxidation states of the same couple, with no
net chemical change, such as $\ce{[Fe(H2O)6]^{3+} + [Fe(H2O)6]^{2+}}$.
\end{definition}

Henry Taube showed the inner-sphere path in 1953 with a reaction chosen so that
the product told its story. Cobalt(III) is inert, cobalt(II) \omterm{def:b3:complex-mechanisms:labile}{labile}; chromium(II)
is \omterm{def:b3:complex-mechanisms:labile}{labile}, chromium(III) inert. When $\ce{[Co(NH3)5Cl]^{2+}}$ is reduced by $\ce{[Cr(H2O)6]^{2+}}$ in
acid, all the chromium(III) formed carries the chloride, as $\ce{[Cr(H2O)5Cl]^{2+}}$, even in
the presence of free chloride: the chloride must have bridged the two metals,
been bound to the inert chromium(III) as soon as the electron passed, and been
released by the \omterm{def:b3:complex-mechanisms:labile}{labile} cobalt(II).

\begin{omfigure}
\begin{tikzpicture}[font=\small]
  \begin{scope}[scale=0.8]
    \pic[transform shape] (a) at (0,0) {omoctahedron};
    \pic[transform shape] (b) at (2.0,0) {omoctahedron};
  \end{scope}
  \node[anchor=north east] at (0.3,-1.0) {$\mathrm{Co^{III}(NH_3)_5}$};
  \node[anchor=north west] at (1.3,-1.0) {$\mathrm{Cr^{II}(H_2O)_5}$};
  \node[omThm, anchor=south] at (0.8,0.15) {Cl};
  \draw[->, thick, omDef] (1.6,1.75) to[bend right=40] node[midway, above] {e$^-$} (0,1.75);
\end{tikzpicture}

{\small Taube's bridged intermediate (schematic): the chloride, still bound to
cobalt(III), enters the coordination sphere of chromium(II) by replacing a water
molecule; the electron crosses through the bridge, and the chloride stays on the
now inert chromium(III).}
\end{omfigure}

An electron jumps in about $10^{-16}$ s, far faster than nuclei move: by the
\omterm{thm:b3:electronic-spectroscopy:franck-condon}{Franck--Condon principle} (\cref{ch:b3:electronic-spectroscopy}), the transfer
can only happen at a nuclear configuration where reactants and products have
the same energy. For the iron self-exchange, the \ce{Fe^{III}}--O bonds are shorter than
the \ce{Fe^{II}}--O bonds; before the electron can move, the two complexes must distort to
a common intermediate geometry, at a cost of energy.

\section{Marcus theory}

\begin{definition}[Reorganisation energy]\label{def:b3:complex-mechanisms:reorganisation}
The \emph{reorganisation energy}\index{reorganisation energy} $\lambda$ of an
electron transfer is the energy needed to bring the reactants, without
transferring the electron, to the nuclear configuration (bond lengths of the
complexes and orientations of the surrounding solvent) of the products. It is the
sum of an inner part, from the bonds, and an outer part, from the solvent.
\end{definition}

\begin{theorem}[Marcus equation]\label{thm:b3:complex-mechanisms:marcus}
If the Gibbs energies of reactants and products are parabolas of the same
curvature in a reaction coordinate $x$ (0 at the reactants' equilibrium, 1 at the
products'), the \omterm{def:b3:rate-theories:activation}{Gibbs energy of activation} of an electron transfer with
standard reaction Gibbs energy $\Delta G^\circ$ and \omterm{def:b3:complex-mechanisms:reorganisation}{reorganisation energy} $\lambda$ is
\[
  \Delta G^{\ddagger} = \frac{(\lambda + \Delta G^\circ)^2}{4\lambda} .
\]
This is the \emph{Marcus equation}\index{Marcus equation}.
\end{theorem}

\begin{proof}
Write $G_R = \lambda x^2$ and $G_P = \lambda(x - 1)^2 + \Delta G^\circ$: the curvature is fixed by $G_R(1) = \lambda$, the
definition of $\lambda$. The transfer occurs where the curves cross: $\lambda x^2 = \lambda x^2 - 2\lambda x + \lambda +
\Delta G^\circ$, so $x^\ast = (\lambda + \Delta G^\circ)/2\lambda$, and $\Delta G^{\ddagger} = G_R(x^\ast) = (\lambda + \Delta G^\circ)^2/4\lambda$.
\end{proof}

\begin{corollary}[Inverted region]\label{cor:b3:complex-mechanisms:inverted}
At fixed $\lambda$, the rate increases as the reaction becomes more exergonic until
$-\Delta G^\circ = \lambda$, where $\Delta G^{\ddagger} = 0$; beyond, it decreases again.
\end{corollary}

\begin{proof}
$\Delta G^{\ddagger}$ is a parabola in $\Delta G^\circ$ with its minimum, zero, at $\Delta G^\circ = -\lambda$; for $\Delta G^\circ < -\lambda$ it grows
again.
\end{proof}

\begin{definition}[Inverted region]\label{def:b3:complex-mechanisms:inverted}
The \emph{inverted region}\index{inverted region} of electron transfer is the
range of driving forces $-\Delta G^\circ > \lambda$ in which the rate decreases as the reaction becomes
more exergonic.
\end{definition}

\begin{omfigure}
\begin{tikzpicture}
\begin{axis}[name=mp, width=0.47\linewidth, height=5.6cm, xmin=-0.6, xmax=2.2, ymin=-160, ymax=120,
  axis lines=left, xlabel={reaction coordinate $x$}, ylabel={$G$ / \unit{kJ.mol^{-1}}},
  tick label style={font=\small}, label style={font=\small}, xtick={0,1,2}]
\addplot[black, thick] table[x=x, y=GR] {figdata/out/bachelor-3-marcus-parabolas.dat};
\addplot[omDef!50, thick] table[x=x, y=P0] {figdata/out/bachelor-3-marcus-parabolas.dat};
\addplot[omDef, thick] table[x=x, y=P50] {figdata/out/bachelor-3-marcus-parabolas.dat};
\addplot[omThm, thick] table[x=x, y=P100] {figdata/out/bachelor-3-marcus-parabolas.dat};
\addplot[omThm!60, thick, dashed] table[x=x, y=P150] {figdata/out/bachelor-3-marcus-parabolas.dat};
\node[font=\scriptsize, anchor=north] at (axis cs:0,-2) {R};
\end{axis}
\begin{axis}[at={(mp.east)}, anchor=west, xshift=1.4cm, width=0.47\linewidth, height=5.6cm,
  xmin=0, xmax=250, ymin=0, ymax=11.5, axis lines=left, xlabel={$-\Delta G^\circ$ / \unit{kJ.mol^{-1}}},
  ylabel={$\log(k/\unit{s^{-1}})$}, tick label style={font=\small}, label style={font=\small}]
\addplot[omDef, very thick] table[x=mdG, y=logk] {figdata/out/bachelor-3-marcus-rate.dat};
\draw[black!40, dashed] (axis cs:100,0) -- (axis cs:100,11);
\node[font=\scriptsize, anchor=south east] at (axis cs:95,1) {normal};
\node[font=\scriptsize, anchor=south west] at (axis cs:105,1) {inverted};
\end{axis}
\end{tikzpicture}

{\small Marcus theory with $\lambda = \qty{100}{kJ/mol}$ (model). Left: the reactant parabola
(black) and product parabolas for $\Delta G^\circ = 0$ and $-50$ (blue), $-100 = -\lambda$ (red) and
\qty{-150}{kJ/mol} (dashed, \omterm{def:b3:complex-mechanisms:inverted}{inverted region}); their crossing
point, the transition state, falls to the bottom of the reactant curve at
$-\Delta G^\circ = \lambda$ and rises again beyond. Right: the logarithm of the rate constant
(model prefactor \qty{e11}{s^{-1}}) passes through a maximum at $-\Delta G^\circ = \lambda$.}
\end{omfigure}

\begin{proposition}[Outer-sphere reorganisation energy]\label{prop:b3:complex-mechanisms:outer-reorganisation}
For two spherical reactants of radii $a_1$, $a_2$ at a distance $d$ in a solvent of
refractive index $n$ and relative permittivity $\varepsilon_s$, the solvent part of $\lambda$ is
proportional to $\bigl(\frac{1}{2a_1} + \frac{1}{2a_2} - \frac1d\bigr)\bigl(\frac{1}{n^2} - \frac{1}{\varepsilon_s}\bigr)$.
\end{proposition}

\admitted

Water, very polar ($\varepsilon_s$ large) but with an ordinary refractive index, gives a large
solvent reorganisation; non-polar solvents a small one. Large reactants
reorganise less: this is why electron-transfer proteins bury their metal sites
inside the protein.

\begin{theorem}[Marcus cross relation]\label{thm:b3:complex-mechanisms:cross-relation}
For a cross reaction $\mathrm{A_{ox} + B_{red} \to A_{red} + B_{ox}}$ with equilibrium constant $K_{12}$, between couples
whose self-exchange rate constants are $k_{11}$ and $k_{22}$,
\[
  k_{12} = \sqrt{k_{11}k_{22}K_{12}f}, \qquad \ln f = \frac{(\ln K_{12})^2}{4\ln(k_{11}k_{22}/Z^2)},
\]
with $Z$ the collision frequency factor; $f \approx 1$ when $K_{12}$ is not too large. This is the
\emph{Marcus cross relation}\index{Marcus cross relation}.
\end{theorem}

\begin{proof}[Partial proof]
Assume the \omterm{def:b3:complex-mechanisms:reorganisation}{reorganisation energy} of the cross reaction is the mean of those of the
self-exchanges, $\lambda_{12} = (\lambda_{11} + \lambda_{22})/2$ (each reactant contributes half of its own
self-exchange reorganisation), and all rate constants $k = Z\eu^{-\Delta G^{\ddagger}/RT}$. Then
$\Delta G_{12}^{\ddagger} = \frac{\lambda_{12}}{4}\bigl(1 + \frac{\Delta G^\circ_{12}}{\lambda_{12}}\bigr)^2 = \frac{\lambda_{12}}{4} + \frac{\Delta G^\circ_{12}}{2} + \frac{(\Delta G^\circ_{12})^2}{4\lambda_{12}}$. The first term gives
$\sqrt{k_{11}k_{22}}$ (since $\Delta G^{\ddagger}_{ii} = \lambda_{ii}/4$), the second $\sqrt{K_{12}}$ (with $\Delta G^\circ_{12} = -RT\ln K_{12}$), the third
the factor $f$.
\end{proof}

\begin{method}[A cross-reaction rate from self-exchange data]\label{met:b3:complex-mechanisms:marcus-estimate}
\begin{enumerate}
  \item Find the self-exchange constants $k_{11}$, $k_{22}$ of the two couples.
  \item Compute $K_{12}$ from the standard potentials: $\ln K_{12} = nF(E^\circ_{\text{oxidant}} - E^\circ_{\text{reductant}})/RT$.
  \item $k_{12} = \sqrt{k_{11}k_{22}K_{12}}$ (with $f$ if $K_{12}$ is very large). Agreement within a factor of a
        few with the measured value indicates an outer-sphere mechanism.
\end{enumerate}
\end{method}

Marcus published the theory in 1956; the \omterm{def:b3:complex-mechanisms:inverted}{inverted region}, at first doubted, was
observed in 1984 by Gerhard Closs and John Miller in molecules where donor and
acceptor are held at a fixed distance by a rigid spacer, so that diffusion
cannot hide it.

\begin{inthelab}[Water exchange by oxygen-17 NMR]
The oxygen-17 resonance of water bound to a paramagnetic ion and that of bulk
water are broadened by the exchange between them. Recording the linewidth of the
bulk water signal against temperature, in a solution of the metal perchlorate
enriched in \ce{^{17}O}, gives the exchange rate constant and, in a high-pressure probe,
its activation volume.
\end{inthelab}

\begin{safety}{\ghs[0.8cm]{GHS05}\ghs[0.8cm]{GHS06}\ghs[0.8cm]{GHS08}}
% ledger: ghs4:K2PtCl4
Potassium tetrachloroplatinate(II): toxic if swallowed, causes serious eye
damage, a respiratory and skin sensitiser. Platinum complexes are handled with
gloves in a fume hood; cisplatin and its analogues are cytotoxic drugs and are
handled as such.
\end{safety}

\begin{history}[Taube and Marcus]
Henry Taube classified complexes as \omterm{def:b3:complex-mechanisms:labile}{labile} or inert in 1952, and with his
chromium--cobalt experiment of 1953 established the inner-sphere mechanism; he
received the 1983 Nobel Prize in Chemistry. Rudolph Marcus derived from 1956 the
theory of electron transfer that bears his name, and received the 1992 Nobel
Prize in Chemistry.
\end{history}

\section{Exercises}

\begin{exercise}[$\star$]\label{exo:b3:complex-mechanisms:1}
Classify as \omterm{def:b3:complex-mechanisms:labile}{labile} or inert, with a reason: $\ce{[Cr(H2O)6]^{3+}}$, $\ce{[Co(NH3)6]^{3+}}$, $\ce{[Cr(H2O)6]^{2+}}$,
$\ce{[Cu(H2O)6]^{2+}}$, $\ce{[Zn(H2O)6]^{2+}}$.
\end{exercise}

\begin{exercise}[$\star$]\label{exo:b3:complex-mechanisms:2}
Show that the D rate law reduces to $v = k_1[\mathrm{ML_5X}]$ at high $[\mathrm Y]$ and to $v = (k_1k_2/k_{-1})
[\mathrm{ML_5X}][\mathrm Y]/[\mathrm X]$ when $k_{-1}[\mathrm X] \gg k_2[\mathrm Y]$.
\end{exercise}

\begin{exercise}[$\star$]\label{exo:b3:complex-mechanisms:3}
Predict the signs of $\Delta V^{\ddagger}$ and $\Delta^{\ddagger}S^\circ$ for a D and an A substitution.
\end{exercise}

\begin{exercise}[$\star$]\label{exo:b3:complex-mechanisms:4}
Give the products of $\ce{[PtCl4]^{2-}}$ with two \ce{NH3}, and of $\ce{[Pt(NH3)4]^{2+}}$ with two \ce{Cl-}.
\end{exercise}

\begin{exercise}[$\star\star$]\label{exo:b3:complex-mechanisms:5}
For a substitution on an aqua ion, $K_{\mathrm{os}} = \qty{2.0}{L.mol^{-1}}$ and $k_i = \qty{3.0e4}{s^{-1}}$ (data of
the exercise). Compute $k_{\mathrm{obs}}$ at $[\mathrm Y] = 0.10$ and \qty{1.0}{mol/L}, and its limit.
\end{exercise}

\begin{exercise}[$\star\star$]\label{exo:b3:complex-mechanisms:6}
A rate constant doubles when the pressure rises from 0.1 to \qty{100}{MPa} at \qty{298}{K}.
Compute $\Delta V^{\ddagger}$ and conclude.
\end{exercise}

\begin{exercise}[$\star\star$]\label{exo:b3:complex-mechanisms:7}
In $\ce{[PtCl3(PEt3)]-}$ the Pt--Cl bond trans to \ce{PEt3} is \qty{2.42}{\AA} and the two others \qty{2.32}{\AA}
(data of the exercise). Interpret, and predict which chloride is replaced first.
\end{exercise}

\begin{exercise}[$\star\star$]\label{exo:b3:complex-mechanisms:8}
List the experimental evidence that would show an electron transfer to be inner
sphere.
\end{exercise}

\begin{exercise}[$\star\star$]\label{exo:b3:complex-mechanisms:9}
For $\lambda = \qty{80}{kJ/mol}$, compute $\Delta G^{\ddagger}$ for $\Delta G^\circ = 0$, $-20$, $-80$ and \qty{-120}{kJ/mol}.
\end{exercise}

\begin{exercise}[$\star\star\star$]\label{exo:b3:complex-mechanisms:10}
Two couples have self-exchange constants $k_{11} = 4.0$ and $k_{22} = \qty{2.0e5}{L.mol^{-1}.s^{-1}}$, and
their cross reaction has $K_{12} = \num{1.0e4}$ (data of the exercise). Estimate $k_{12}$ ($f = 1$).
\end{exercise}

\begin{exercise}[$\star\star\star$]\label{exo:b3:complex-mechanisms:11}
Using the $\sigma$-only model, compute the ligand-field activation energies of
high-spin $d^5$, $d^7$ and $d^8$ ions and rank \ce{Mn^{2+}}, \ce{Co^{2+}} and \ce{Ni^{2+}} by expected
water-exchange rate.
\end{exercise}

\begin{exercise}[$\star\star\star$]\label{exo:b3:complex-mechanisms:12}
Why is the \omterm{def:b3:complex-mechanisms:inverted}{inverted region} hard to observe for reactions between ions that
diffuse together in solution, and how did rigid donor--acceptor molecules reveal
it?
\end{exercise}

\section{Problem: Making Cisplatin, Not Transplatin}

\begin{problem}[{Weekend problem --- making cisplatin, not transplatin: the
square-planar platinum(II) ion, a synthesis planned with the trans effect, the
aquation of the drug, and why it is activated inside cells}]
\label{pb:b3:complex-mechanisms:1}
% ledger: ghs4:K2PtCl4
Cisplatin, $cis$-$\ce{[PtCl2(NH3)2]}$, is made from $\ce{K2[PtCl4]}$. Data of the problem: at
\qty{37}{\celsius} the first aquation, $\ce{[PtCl2(NH3)2] + H2O -> [PtCl(H2O)(NH3)2]+ + Cl-}$, has the
rate constant $k = \qty{9.0e-5}{s^{-1}}$ and the equilibrium constant $K = \qty{3.0e-3}{mol/L}$; the
chloride concentration is about \qty{100}{mM} in blood plasma and \qty{4}{mM} inside cells.

\medskip\noindent\textbf{Part I --- Platinum(II).}
\begin{enumerate}
  \item Give the $d$-electron count of \ce{Pt^{2+}}.
  \item Why are its complexes square planar and diamagnetic?
  \item Why are their substitutions associative?
  \item Write the two-term rate law and explain its two terms.
  \item Is the drug \omterm{def:b3:complex-mechanisms:labile}{labile} or inert on the time scale of a day?
\end{enumerate}

\medskip\noindent\textbf{Part II --- The synthesis.}
\begin{enumerate}[resume]
  \item What does $\ce{K2[PtCl4]}$ give with two \ce{NH3}? Why?
  \item That direct route gives impure product. Dhara's route first converts
        $\ce{[PtCl4]^{2-}}$ into $\ce{[PtI4]^{2-}}$ with excess KI. Why is iodide useful?
  \item Adding \ce{NH3} gives $cis$-$\ce{[PtI2(NH3)2]}$. Explain the stereochemistry.
  \item The iodides are then removed with silver nitrate in water. What forms?
  \item Adding KCl gives cisplatin. Why does the configuration not change?
  \item How would you make transplatin instead?
  \item Why is transplatin of no use as a drug, although it reacts too?
\end{enumerate}

\medskip\noindent\textbf{Part III --- Aquation.}
\begin{enumerate}[resume]
  \item Compute the half-life of the first aquation.
  \item Compute the fraction aquated after \qty{2.0}{h} in a chloride-free solution.
  \item Why is the aqua complex the reactive form towards DNA?
  \item Is the aquation dissociative or associative? What would its activation
        volume be?
  \item Why is a drug solution for injection made up in saline?
\end{enumerate}

\medskip\noindent\textbf{Part IV --- Plasma and cell.}
\begin{enumerate}[resume]
  \item Write the equilibrium fraction of the aqua complex as a function of
        $[\ce{Cl-}]$.
  \item Compute it in plasma.
  \item Compute it inside a cell.
  \item Compute the ratio.
  \item Why does the drug therefore act mostly inside cells?
  \item What other ligands inside a cell can bind platinum and deactivate the
        drug?
  \item Why do the two ammine ligands stay bound throughout?
  \item State the result: the ratio of the aqua-complex fraction inside a cell to
        that in plasma.
\end{enumerate}
\end{problem}
